%% bare_conf.tex
%% V1.2
%% 2002/11/18
%% by Michael Shell
%% mshell@ece.gatech.edu
%%
%% NOTE: This text file uses MS Windows line feed conventions. When (human)
%% reading this file on other platforms, you may have to use a text
%% editor that can handle lines terminated by the MS Windows line feed
%% characters (0x0D 0x0A).
%%
%% This is a skeleton file demonstrating the use of IEEEtran.cls
%% (requires IEEEtran.cls version 1.6b or later) with an IEEE conference paper.
%%
%% Support sites:
%% http://www.ieee.org
%% and/or
%% http://www.ctan.org/tex-archive/macros/latex/contrib/supported/IEEEtran/
%%
%% This code is offered as-is - no warranty - user assumes all risk.
%% Free to use, distribute and modify.

% *** Authors should verify (and, if needed, correct) their LaTeX system  ***
% *** with the testflow diagnostic prior to trusting their LaTeX platform ***
% *** with production work. IEEE's font choices can trigger bugs that do  ***
% *** not appear when using other class files.                            ***
% Testflow can be obtained at:
% http://www.ctan.org/tex-archive/macros/latex/contrib/supported/IEEEtran/testflow


% Note that the a4paper option is mainly intended so that authors in
% countries using A4 can easily print to A4 and see how their papers will
% look in print. Authors are encouraged to use U.S. letter paper when
% submitting to IEEE. Use the testflow package mentioned above to verify
% correct handling of both paper sizes by the author's LaTeX system.
%
% Also note that the "draftcls" or "draftclsnofoot", not "draft", option
% should be used if it is desired that the figures are to be displayed in
% draft mode.
%
% This paper can be formatted using the peerreviewca
% (instead of conference) mode.
% If the IEEEtran.cls has not been installed into the LaTeX system files,
% manually specify the path to it:
% \documentclass[conference]{IEEEtran}
% To use A4 paper, add a4paper option as in
% \documentclass[conference,a4paper]{IEEEtran}


% setup page to suit conference specification using fancyhdr
\documentclass[conference,letterpaper]{IEEEtran}
\usepackage{fancyhdr}
\setlength{\paperwidth}{215.9mm}
\setlength{\hoffset}{-9.7mm}
\setlength{\oddsidemargin}{0mm}
\setlength{\textwidth}{184.3mm}
\setlength{\columnsep}{6.3mm}
\setlength{\marginparsep}{0mm}
\setlength{\marginparwidth}{0mm}

\setlength{\paperheight}{279.4mm}
\setlength{\voffset}{-7.4mm}
\setlength{\topmargin}{0mm}
\setlength{\headheight}{0mm}
\setlength{\headsep}{0mm}
\setlength{\topskip}{0mm}
\setlength{\textheight}{235.2mm}
\setlength{\footskip}{12.4mm}

\setlength{\parindent}{1pc}


% some very useful LaTeX packages include:

%\usepackage{cite}      % Written by Donald Arseneau
                        % V1.6 and later of IEEEtran pre-defines the format
                        % of the cite.sty package \cite{} output to follow
                        % that of IEEE. Loading the cite package will
                        % result in citation numbers being automatically
                        % sorted and properly "ranged". i.e.,
                        % [1], [9], [2], [7], [5], [6]
                        % (without using cite.sty)
                        % will become:
                        % [1], [2], [5]--[7], [9] (using cite.sty)
                        % cite.sty's \cite will automatically add leading
                        % space, if needed. Use cite.sty's noadjust option
                        % (cite.sty V3.8 and later) if you want to turn this
                        % off. cite.sty is already installed on most LaTeX
                        % systems. The latest version can be obtained at:
                        % http://www.ctan.org/tex-archive/macros/latex/contrib/supported/cite/

%\usepackage{graphicx}  % Written by David Carlisle and Sebastian Rahtz
                        % Required if you want graphics, photos, etc.
                        % graphicx.sty is already installed on most LaTeX
                        % systems. The latest version and documentation can
                        % be obtained at:
                        % http://www.ctan.org/tex-archive/macros/latex/required/graphics/
                        % Another good source of documentation is "Using
                        % Imported Graphics in LaTeX2e" by Keith Reckdahl
                        % which can be found as esplatex.ps and epslatex.pdf
                        % at: http://www.ctan.org/tex-archive/info/
                        % NOTE: for dual use with latex and pdflatex, instead load graphicx like:
                        %\ifx\pdfoutput\undefined
                        %\usepackage{graphicx}
                        %\else
                        %\usepackage[pdftex]{graphicx}
                        %\fi
                        %
                        % However, be warned that pdflatex will require graphics to be in PDF
                        % (not EPS) format and will preclude the use of PostScript based LaTeX
                        % packages such as psfrag.sty and pstricks.sty. IEEE conferences typically
                        % allow PDF graphics (and hence pdfLaTeX). However, IEEE journals do not
                        % (yet) allow image formats other than EPS or TIFF. Therefore, authors of
                        % journal papers should use traditional LaTeX with EPS graphics.
                        %
                        % The path(s) to the graphics files can also be declared: e.g.,
                        % \graphicspath{{../eps/}{../ps/}}
                        % if the graphics files are not located in the same directory as the
                        % .tex file. This can be done in each branch of the conditional above
                        % (after graphicx is loaded) to handle the EPS and PDF cases separately.
                        % In this way, full path information will not have to be specified in
                        % each \includegraphics command.
                        %
                        % Note that, when switching from latex to pdflatex and vice-versa, the new
                        % compiler will have to be run twice to clear some warnings.

%\usepackage{psfrag}    % Written by Craig Barratt, Michael C. Grant,
                        % and David Carlisle
                        % This package allows you to substitute LaTeX
                        % commands for text in imported EPS graphic files.
                        % In this way, LaTeX symbols can be placed into
                        % graphics that have been generated by other
                        % applications. You must use latex->dvips->ps2pdf
                        % workflow (not direct pdf output from pdflatex) if
                        % you wish to use this capability because it works
                        % via some PostScript tricks. Alternatively, the
                        % graphics could be processed as separate files via
                        % psfrag and dvips, then converted to PDF for
                        % inclusion in the main file which uses pdflatex.
                        % Docs are in "The PSfrag System" by Michael C. Grant
                        % and David Carlisle. There is also some information
                        % about using psfrag in "Using Imported Graphics in
                        % LaTeX2e" by Keith Reckdahl which documents the
                        % graphicx package (see above). The psfrag package
                        % and documentation can be obtained at:
                        % http://www.ctan.org/tex-archive/macros/latex/contrib/supported/psfrag/

%\usepackage{subfigure} % Written by Steven Douglas Cochran
                        % This package makes it easy to put subfigures
                        % in your figures. i.e., "figure 1a and 1b"
                        % Docs are in "Using Imported Graphics in LaTeX2e"
                        % by Keith Reckdahl which also documents the graphicx
                        % package (see above). subfigure.sty is already
                        % installed on most LaTeX systems. The latest version
                        % and documentation can be obtained at:
                        % http://www.ctan.org/tex-archive/macros/latex/contrib/supported/subfigure/

%\usepackage{url}       % Written by Donald Arseneau
                        % Provides better support for handling and breaking
                        % URLs. url.sty is already installed on most LaTeX
                        % systems. The latest version can be obtained at:
                        % http://www.ctan.org/tex-archive/macros/latex/contrib/other/misc/
                        % Read the url.sty source comments for usage information.

%\usepackage{stfloats}  % Written by Sigitas Tolusis
                        % Gives LaTeX2e the ability to do double column
                        % floats at the bottom of the page as well as the top.
                        % (e.g., "\begin{figure*}[!b]" is not normally
                        % possible in LaTeX2e). This is an invasive package
                        % which rewrites many portions of the LaTeX2e output
                        % routines. It may not work with other packages that
                        % modify the LaTeX2e output routine and/or with other
                        % versions of LaTeX. The latest version and
                        % documentation can be obtained at:
                        % http://www.ctan.org/tex-archive/macros/latex/contrib/supported/sttools/
                        % Documentation is contained in the stfloats.sty
                        % comments as well as in the presfull.pdf file.
                        % Do not use the stfloats baselinefloat ability as
                        % IEEE does not allow \baselineskip to stretch.
                        % Authors submitting work to the IEEE should note
                        % that IEEE rarely uses double column equations and
                        % that authors should try to avoid such use.
                        % Do not be tempted to use the cuted.sty or
                        % midfloat.sty package (by the same author) as IEEE
                        % does not format its papers in such ways.

%\usepackage{amsmath}   % From the American Mathematical Society
                        % A popular package that provides many helpful commands
                        % for dealing with mathematics. Note that the AMSmath
                        % package sets \interdisplaylinepenalty to 10000 thus
                        % preventing page breaks from occurring within multiline
                        % equations. Use:
                        %\interdisplaylinepenalty=2500
                        % after loading amsmath to restore such page breaks
                        % as IEEEtran.cls normally does. amsmath.sty is already
                        % installed on most LaTeX systems. The latest version
                        % and documentation can be obtained at:
                        % http://www.ctan.org/tex-archive/macros/latex/required/amslatex/math/


% Other popular packages for formatting tables and equations include:

%\usepackage{array}
% Frank Mittelbach's and David Carlisle's array.sty which improves the
% LaTeX2e array and tabular environments to provide better appearances and
% additional user controls. array.sty is already installed on most systems.
% The latest version and documentation can be obtained at:
% http://www.ctan.org/tex-archive/macros/latex/required/tools/

% Mark Wooding's extremely powerful MDW tools, especially mdwmath.sty and
% mdwtab.sty which are used to format equations and tables, respectively.
% The MDWtools set is already installed on most LaTeX systems. The lastest
% version and documentation is available at:
% http://www.ctan.org/tex-archive/macros/latex/contrib/supported/mdwtools/


% V1.6 of IEEEtran contains the IEEEeqnarray family of commands that can
% be used to generate multiline equations as well as matrices, tables, etc.


% Also of notable interest:

% Scott Pakin's eqparbox package for creating (automatically sized) equal
% width boxes. Available:
% http://www.ctan.org/tex-archive/macros/latex/contrib/supported/eqparbox/


% Notes on hyperref:
% IEEEtran.cls attempts to be compliant with the hyperref package, written
% by Heiko Oberdiek and Sebastian Rahtz, which provides hyperlinks within
% a document as well as an index for PDF files (produced via pdflatex).
% However, it is a tad difficult to properly interface LaTeX classes and
% packages with this (necessarily) complex and invasive package. It is
% recommended that hyperref not be used for work that is to be submitted
% to the IEEE. Users who wish to use hyperref *must* ensure that their
% hyperref version is 6.72u or later *and* IEEEtran.cls is version 1.6b
% or later. The latest version of hyperref can be obtained at:
%
% http://www.ctan.org/tex-archive/macros/latex/contrib/supported/hyperref/
%
% Also, be aware that cite.sty (as of version 3.9, 11/2001) and hyperref.sty
% (as of version 6.72t, 2002/07/25) do not work optimally together.
% To mediate the differences between these two packages, IEEEtran.cls, as
% of v1.6b, predefines a command that fools hyperref into thinking that
% the natbib package is being used - causing it not to modify the existing
% citation commands, and allowing cite.sty to operate as normal. However,
% as a result, citation numbers will not be hyperlinked. Another side effect
% of this approach is that the natbib.sty package will not properly load
% under IEEEtran.cls. However, current versions of natbib are not capable
% of compressing and sorting citation numbers in IEEE's style - so this
% should not be an issue. If, for some strange reason, the user wants to
% load natbib.sty under IEEEtran.cls, the following code must be placed
% before natbib.sty can be loaded:
%
% \makeatletter
% \let\NAT@parse\undefined
% \makeatother
%
% Hyperref should be loaded differently depending on whether pdflatex
% or traditional latex is being used:
%
%\ifx\pdfoutput\undefined
%\usepackage[hypertex]{hyperref}
%\else
%\usepackage[pdftex,hypertexnames=false]{hyperref}
%\fi
%
% Pdflatex produces superior hyperref results and is the recommended
% compiler for such use.

\usepackage{fancyhdr}


% *** Do not adjust lengths that control margins, column widths, etc. ***
% *** Do not use packages that alter fonts (such as pslatex).         ***
% There should be no need to do such things with IEEEtran.cls V1.6 and later.


% correct bad hyphenation here
\hyphenation{op-tical net-works semi-conduc-tor IEEEtran}


\begin{document}
% paper title
\title{Paper Title}


% author names and affiliations
% use a multiple column layout for up to three different
% affiliations
%\author{\IEEEauthorblockN{Authors Name/s per 1st Affiliation}
%\IEEEauthorblockA{line 1 (of Affiliation): dept. name of organization\\
%line 2: name of organization, acronyms acceptable\\
%line 3: City, Country\\
%line 4: e-mail address if desired}
%\and
%\IEEEauthorblockN{Authors Name/s per 2nd Affiliation}
%\IEEEauthorblockA{line 1 (of Affiliation): dept. name of organization\\
%line 2: name of organization, acronyms acceptable\\
%line 3: City, Country\\
%line 4: e-mail address if desired}}


% avoiding spaces at the end of the author lines is not a problem with
% conference papers because we don't use \thanks or \IEEEmembership

\author{\authorblockN{Benjamin M. Rodriguez}
\authorblockA{Johns Hopkins University Applied Physics Laboratory\\
Laurel, Maryland 20723\\ Email: benjamin.rodriguez@jhuapl.edu}}


% for over three affiliations, or if they all won't fit within the width
% of the page, use this alternative format:
%
%\author{\authorblockN{Michael Shell\authorrefmark{1},
%Homer Simpson\authorrefmark{2},
%James Kirk\authorrefmark{3},
%Montgomery Scott\authorrefmark{3} and
%Eldon Tyrell\authorrefmark{4}}
%\authorblockA{\authorrefmark{1}School of Electrical and Computer Engineering\\
%Georgia Institute of Technology,
%Atlanta, Georgia 30332--0250\\ Email: mshell@ece.gatech.edu}
%\authorblockA{\authorrefmark{2}Twentieth Century Fox, Springfield, USA\\
%Email: homer@thesimpsons.com}
%\authorblockA{\authorrefmark{3}Starfleet Academy, San Francisco, California 96678-2391\\
%Telephone: (800) 555--1212, Fax: (888) 555--1212}
%\authorblockA{\authorrefmark{4}Tyrell Inc., 123 Replicant Street, Los Angeles, California 90210--4321}}


% use only for invited papers
%\specialpapernotice{(Invited Paper)}

% make the title area
\maketitle
\thispagestyle{plain}


% insert page header and footer here for IEEE PDF Compliant
%\fancypagestyle{plain}{
%\fancyhf{}	% clear all header and footer fields
%\fancyfoot[L]{978-1-4244-2794-9/09/\$25.00~\copyright~2009 IEEE}
%\fancyfoot[C]{}
%\fancyfoot[R]{}
%\renewcommand{\headrulewidth}{0pt}
%\renewcommand{\footrulewidth}{0pt}
%}


\pagestyle{fancy}{
\fancyhf{}
\fancyfoot[R]{}}
\renewcommand{\headrulewidth}{0pt}
\renewcommand{\footrulewidth}{0pt}



\begin{abstract}

The information in this template is a work in progress from multiple sources and by no means is one author being credited for the work as original. The intention of this template is for research and educational purposes.\\

The abstract should contain the following:
    \begin{itemize}
        \item One or two sentences:  Motivation \\
        \item One sentence:   Topic/Objective What I am doing \\
        \item One or two sentences:  Supporting my work \\
        \item One or two sentences:  Comparison \\
        \item One sentence:   Conclusion \\
    \end{itemize}
\end{abstract}


\begin{IEEEkeywords}
Keywords goes here.
\end{IEEEkeywords}


% For peer review papers, you can put extra information on the cover
% page as needed:
% \begin{center} \bfseries EDICS Category: 3-BBND \end{center}
%
% for peerreview papers, inserts a page break and creates the second title.
% Will be ignored for other modes.
\IEEEpeerreviewmaketitle


\section{Introduction}
% no \PARstart
The introduction section should contain the following paragraphs:
    \begin{itemize}
        \item Paragraph 1 motivation: broadly, what is problem area, why is it important?

        \item Paragraph 2 narrow down: what is the problem I am specifically consider

        \item Paragraph 3 elevator pitch : "In the paper, we \dots": most crucial paragraph

        \item Paragraph 4 contribution: how different/better/relates to other work

        \item Paragraph 5 paper layout: "The remainder of this paper is structured as follows. In Section 2, we introduce \dots Section 3 describes ... Finding and analysis are provided in Section 4. Finally, the conclusion and future work are described in Section 5." [Note that Section is capitalized. Also, vary your expression between "section" being the subject of the sentence, as in "Section 2 discusses ..." and "In Section, we discuss ...".]
    \end{itemize}
%\hfill mds
%
%\hfill November 18, 2002



\section{Literature Review}
In this section provide enough information in a technical document to allow the reader to understand the specific problem being addressed and to provide a context for the document.
\subsection{Subsection Heading Here}
The background information may include:
    \begin{itemize}
	   \item historical summary of the problem being addressed;
	   \item a brief summary of previous work on the topic, including, if appropriate, relevant theory;
       \item related work that is being directly used or being benchmarked should be described.
    \end{itemize}
\subsubsection{Hint}
In the case of a conference, make sure to cite the work of the PC co-chairs and as many other PC members as are remotely plausible, as well as from anything relevant from the previous two proceedings. In the case of a journal or magazine, cite anything relevant from last 2-3 years or so volumes.

\subsubsection{Short Documents}
In short documents, include background information in the introduction. In longer documents, however, putting some or all of the background information in a separate section with a heading may be more effective. Long and fairly complex reports, especially experimental reports where the purpose of the document is to verify, evaluate, illustrate, or apply one or more theories, often include a separate theory section.

\subsubsection{Long Reports and Articles}
In long and fairly complex reports and articles, especially theoretical and experimental reports where the purpose of the document is to apply, verify, or illustrate one or more theories, include a separate section presenting relevant theoretical formulae and the techniques by which any experimental results are predicted. When introducing equations, be sure to define all symbols used in them.



\section{Methodology}

In this section the methodology used in the paper should be detailed. Details in this section should include the input, algorithm, mathematical description, output and other relevant information.

The proposed model is described. Invariably there are assumptions made. State the model assumptions clearly. Do the assumptions make sense? Sometimes it will be necessary to introduce assumptions to make the problem mathematically tractable, but they should at least reflect some real-world situations. Thus while the proposed assumptions may not hold in general, there should at least be some instances where they will hold and hence the model and conclusions drawn will apply. Use figures to help explain the model.

\subsection{Body of the Paper}
    \begin{enumerate}
        \item problem
        \item approach, architecture
        \item results
    \end{enumerate}
The body should contain sufficient motivation, with at least one example scenario, preferably two, with illustrating figures, followed by a crisp generic problem statement model, i.e., functionality, particularly emphasizing "new" functionality. The paper may or may not include formalisms. General evaluations of your algorithm or architecture, e.g., material proving that the algorithm is O(log N), go here, not in the evaluation section.
\subsubsection{Architecture}
Architecture of proposed system(s) to achieve this model should be more generic than your own peculiar implementation. Always include at least one figure.
\subsubsection{Realization}
The realization contains actual implementation details when implementing architecture isn't totally straightforward. Mention briefly implementation language, platform, location, dependencies on other packages and minimum resource usage if pertinent.
\subsubsection{Evaluation (Findings and Analysis)}
In the evaluation, how does it really work in practice? Provide real or simulated performance metrics, end-user studies, mention external technology adaptors, if any, etc.

\section{Findings and Analysis}

This section will include any input parameters, benchmark description, performance measures, tables, graphs, comparisons and analysis. Based on the model, numerical results will be generated. These results should be presented in such a way as to facilitate the readers' understanding. Usually, they will be presented in the form of figures or tables. The parameter values chosen should make sense. They should preferably be taken from systems in the real world. If that is not possible, they may correspond to those values for which published results are available so that one can compare these results with existing ones. All the results should be interpreted. Try to explain why the curves look the way they do. Is it because of the assumptions, or because the system behaves that way? In most cases, a simulation model is required to validate the system model with the assumptions. In that case, it is important to not only show the average values of the simulation results, but also confidence intervals. A number of books on modeling techniques will show how confidence intervals may be obtained. In addition, it is very important to explain what is being simulated. I once attended a thesis defense in which the candidate claimed that his analytical results matched the simulation results perfectly. I then found that he had basically simulated the analytical model, using the same assumptions! Details on the simulation time, the computer, and the language used in the simulation should also be included.

% Reminder: the "draftcls" or "draftclsnofoot", not "draft", class option
% should be used if it is desired that the figures are to be displayed while
% in draft mode.


% An example of a floating figure using the graphicx package.
% Note that \label must occur AFTER (or within) \caption.
% For figures, \caption should occur after the \includegraphics.
%
% \begin{figure}
% \centering
% \includegraphics[width=2.5in]{myfigure}
% where an .eps filename suffix will be assumed under latex,
% and a .pdf suffix will be assumed for pdflatex
% \caption{Simulation Results}
% \label{fig_sim}
% \end{figure}


% An example of a double column floating figure using two subfigures.
% (The subfigure.sty package must be loaded for this to work.)
% The subfigure \label commands are set within each subfigure command, the
% \label for the overall fgure must come after \caption.
% \hfil must be used as a separator to get equal spacing
%
% \begin{figure*}
% \centerline{\subfigure[Case I]{\includegraphics[width=2.5in]{subfigcase1}
% where an .eps filename suffix will be assumed under latex,
% and a .pdf suffix will be assumed for pdflatex
% \label{fig_first_case}}
% \hfil
% \subfigure[Case II]{\includegraphics[width=2.5in]{subfigcase2}
% where an .eps filename suffix will be assumed under latex,
% and a .pdf suffix will be assumed for pdflatex
% \label{fig_second_case}}}
% \caption{Simulation results}
% \label{fig_sim}
% \end{figure*}


% An example of a floating table. Note that, for IEEE style tables, the
% \caption command should come BEFORE the table. Table text will default to
% \footnotesize as IEEE normally uses this smaller font for tables.
% The \label must come after \caption as always.
%
% \begin{table}
% increase table row spacing, adjust to taste
% \renewcommand{\arraystretch}{1.3}
% \caption{An Example of a Table}
% \label{table_example}
% \begin{center}
% Some packages, such as MDW tools, offer better commands for making tables
% than the plain LaTeX2e tabular which is used here.
% \begin{tabular}{|c||c|}
% \hline
% One & Two\\
% \hline
% Three & Four\\
% \hline
% \end{tabular}
% \end{center}
% \end{table}


\section{Conclusion}
The conclusion summarizes what have been done and concluded based on the results. A description of future research should also be included.
    \begin{itemize}
	   \item One sentence '  your contribution (proposed method)
	   \item Two or more sentences '  comparison results (The advantages of the presented technique are)
	   \item One or two sentences ' conclude with remarks that your contribution is accountable
	   \item One or two sentences ' stating future work
    \end{itemize}
Easiest thing to do is that just reword/rewrite your abstract



% conference papers do not normally have an appendix


% use section* for acknowledgement
\section*{Acknowledgment}
% optional entry into table of contents (if used)
% \addcontentsline{toc}{section}{Acknowledgment}
The authors would like to thank...


% trigger a \newpage just before the given reference
% number - used to balance the columns on the last page
% adjust value as needed - may need to be readjusted if
% the document is modified later
% \IEEEtriggeratref{8}
% The "triggered" command can be changed if desired:
% \IEEEtriggercmd{\enlargethispage{-5in}}


% references section
% NOTE: BibTeX documentation can be easily obtained at:
% http://www.ctan.org/tex-archive/biblio/bibtex/contrib/doc/


% can use a bibliography generated by BibTeX as a .bbl file
% standard IEEE bibliography style from:
% http://www.ctan.org/tex-archive/macros/latex/contrib/supported/IEEEtran/bibtex
% \bibliographystyle{IEEEtran.bst}
% argument is your BibTeX string definitions and bibliography database(s)
% \bibliography{IEEEabrv,../bib/paper}
%
% <OR> manually copy in the resultant .bbl file
% set second argument of \begin to the number of references
% (used to reserve space for the reference number labels box)
\begin{thebibliography}{1}

\bibitem{IEEEhowto:kopka}
H.~Kopka and P.~W. Daly, \emph{A Guide to {\LaTeX}}, 3rd~ed.\hskip 1em plus
  0.5em minus 0.4em\relax Harlow, England: Addison-Wesley, 1999.
\bibitem{IEEEhowto:KaiserEtAl}  
  G.~Kaiser, C.~Partridge, S.~Roy, E.~Siegel, S.~Stolfo, L.~Trevisan, Y.~Yemini and E.~Zadok, \emph{Writing Technical Articles}, Retrived Jan 5, 2008 from http://www.cs.columbia.edu/~hgs/etc/writing-style.html
  
\end{thebibliography}


% that's all folks
\end{document} 